\scanpage{254}
% archive sheet number: 253
\unsure{constant|constante|constants}. En effet (\unsure{cf|réf} ?) $H^i(X_n,\mathbb Z)=\deleted{\unsure{H|N}}H^i(X_0,\mathbb Z)$ est un groupe de type fini, et son quotient par $\operatorname{Ker}\alpha^i$ est sans torsion [car \unsure{isom} à un sous-groupe de $H^i(X_n,\mathcal O_{X_n})$, qui est un vectoriel sur $\mathbb C$], \deleted{\unsure{c'est|d'}} il suffit donc de noter que pour $n$ grand, \unsure{le|les|la} \unsure{rang|rangs} desdits quotients \unsure{reste|vont} \unsure{constant|constants|constante}.

On obtient ainsi, compte tenu de 1.3 (ii), et du fait qu'un \unsure{ext.|cat.} d'un syst. M.L. par un syst. M.L. est ML

\uline{\deleted{\unsure{Corollaire}}\inserted{Proposition}} 1.5\quad Sous les conditions de 1.3,
(i) Le syst. projectif $(H^i(X_n,\mathcal O_{X_n}^{\times}))_n$ satisfait ML

(ii) On a une suite exacte
\[
H^i(X_0,\mathbb Z)\xrightarrow{\alpha_\infty^{(i)}}
\deleted{\unsure{H|R}}R^if_*(\mathcal O_X)_y^{\wedge}
\xrightarrow{\beta_\infty^{(i)}}\underset{\leftarrow}{\lim}H^i(X_n,\mathcal O_{X_n}^{\times})
\xrightarrow{\unsure{\partial|\delta}_\infty^{(i)}} H^{i+1}(X_0,\mathbb Z)
\xrightarrow{\alpha_\infty^{(i+1)}} R^{i+1}f_*(\mathcal O_X)_y^{\wedge}
\]
\deleted{\illegible{1}}

Notons qu'on a une suite exacte déduite de la suite exacte
\[
0\longrightarrow \deleted{\unsure{\mathcal O|0}}\mathbb Z_X\longrightarrow \uline{\mathcal O}_X\longrightarrow \mathcal O_X^{\times}\longrightarrow 0
\]
par application de $(R^if_*)_y$ et compte tenu de 1.2, les \inserted{fibres en $y$ des} termes $R^if_*(\mathbb Z_X)_{\unsure{y|x}}$ \deleted{se} s'identifient à $H^i(X_0,\mathbb Z)$. On obtient ainsi une
